\section{Coordinación y Efecto Látigo (Canvas)}

\begin{tcolorbox}[colback=cyan!5!white,colframe=cyan!75!black,title=Mapa Mental Rápido (Cópialo)]
\textbf{Efecto Látigo (Bullwhip):} Variación amplificada aguas arriba. Un pequeño cambio en el cliente final genera un caos de pedidos en la fábrica.\\
\textbf{La Solución Maestra:} Compartir el POS (Point of Sale) para que todos vean la demanda real, no los pedidos distorsionados del eslabón anterior.
\end{tcolorbox}

\subsection{Ejercicio Clave de Clase: Doble Marginalización}

\begin{definicion}{Qué está midiendo este ejercicio}
Este ejercicio compara una cadena \textbf{descentralizada} contra una cadena \textbf{centralizada}. En la cadena descentralizada, el proveedor y el minorista maximizan sus propias utilidades por separado. En la cadena centralizada, se maximiza la utilidad total de la cadena completa.

La pérdida de eficiencia aparece porque ambos agentes agregan margen: primero el proveedor sube el precio mayorista \(w\), y luego el minorista sube el precio final \(p\). A esto se le llama \textbf{doble marginalización}.
\end{definicion}

\begin{formulas}{Plantilla del modelo fabricante--minorista}
Demanda inversa:
\[
p=a-bq.
\]
Costo unitario del proveedor:
\[
c.
\]

En la cadena descentralizada, el proveedor elige \(w\) y el minorista elige \(q\). Dado \(w\), el minorista maximiza:
\[
\Pi_r(q)=q(p-w)=q(a-bq-w).
\]
Derivando:
\[
\frac{d\Pi_r}{dq}=a-w-2bq=0
\quad\Rightarrow\quad
q(w)=\frac{a-w}{2b}.
\]
El proveedor anticipa esta respuesta y maximiza:
\[
\Pi_s(w)=(w-c)q(w)
=(w-c)\frac{a-w}{2b}.
\]
Derivando:
\[
\frac{d\Pi_s}{dw}=0
\quad\Rightarrow\quad
w^*=\frac{a+c}{2}.
\]
Por lo tanto:
\[
q_D=\frac{a-c}{4b},
\qquad
p_D=\frac{3a+c}{4}.
\]
Las utilidades descentralizadas son:
\[
\Pi_s^D=\frac{(a-c)^2}{8b},
\qquad
\Pi_r^D=\frac{(a-c)^2}{16b},
\qquad
\Pi_{SC}^D=\frac{3(a-c)^2}{16b}.
\]

En la cadena centralizada, se maximiza la utilidad total:
\[
\Pi_{SC}(q)=q(p-c)=q(a-bq-c).
\]
Derivando:
\[
\frac{d\Pi_{SC}}{dq}=a-c-2bq=0
\quad\Rightarrow\quad
q_C=\frac{a-c}{2b}.
\]
Luego:
\[
p_C=\frac{a+c}{2},
\qquad
\Pi_{SC}^C=\frac{(a-c)^2}{4b}.
\]
\end{formulas}

\begin{ejercicio}{Caso de clase: \(Q=100-P\), \(c=1\)}
El enunciado de clase usa:
\[
Q=100-P
\quad\Leftrightarrow\quad
p=100-q,
\qquad
a=100,\quad b=1,\quad c=1.
\]

\textbf{Cadena descentralizada:}
\[
w_D=\frac{100+1}{2}=50.5,
\]
\[
q_D=\frac{100-1}{4}=24.75,
\qquad
p_D=\frac{3(100)+1}{4}=75.25.
\]
Utilidades:
\[
\Pi_s^D=(50.5-1)(24.75)=1225.125,
\]
\[
\Pi_r^D=(75.25-50.5)(24.75)=612.5625,
\]
\[
\Pi_{SC}^D=1837.6875.
\]

\textbf{Cadena centralizada:}
\[
q_C=\frac{100-1}{2}=49.5,
\qquad
p_C=\frac{100+1}{2}=50.5,
\]
\[
\Pi_{SC}^C=(50.5-1)(49.5)=2450.25.
\]

Comparación:
\[
q_C=49.5>24.75=q_D,
\qquad
p_C=50.5<75.25=p_D,
\]
\[
\Pi_{SC}^C=2450.25>1837.69=\Pi_{SC}^D.
\]
La cadena descentralizada vende menos, cobra más caro y genera menor utilidad total.
\end{ejercicio}

\begin{tipprueba}{Cómo contestarlo rápido}
Si aparece este ejercicio, no partas inventando. Usa esta secuencia:
\[
\text{minorista: } q(w)
\quad\rightarrow\quad
\text{proveedor: } w^*
\quad\rightarrow\quad
\text{descentralizado}
\]
\[
\quad\rightarrow\quad
\text{centralizado}
\quad\rightarrow\quad
\text{comparar}.
\]
La frase final debe decir: ``La falta de coordinación produce doble marginalización; por eso la cadena descentralizada vende menos que el óptimo integrado''.
\end{tipprueba}

\begin{formulas}{Contrato de compartir utilidades}
Una forma de coordinar es un contrato de compartir ingresos/utilidades. Si el minorista conserva una fracción \(\alpha\) de los ingresos y paga \(w\) por unidad, su utilidad es:
\[
\Pi_r(q)=\alpha R(q)-wq.
\]
Para que el minorista elija la cantidad centralizada, debe cumplirse:
\[
R'(q)=c.
\]
Pero el minorista elige según:
\[
\alpha R'(q)-w=0
\quad\Rightarrow\quad
R'(q)=\frac{w}{\alpha}.
\]
Entonces el contrato coordina si:
\[
\frac{w}{\alpha}=c
\quad\Rightarrow\quad
w=\alpha c.
\]
Con \(w=\alpha c\), el minorista recibe:
\[
\Pi_r=\alpha \Pi_{SC}^C,
\]
y el proveedor recibe:
\[
\Pi_s=(1-\alpha)\Pi_{SC}^C.
\]
En el ejemplo:
\[
\Pi_{SC}^C=2450.25.
\]
Para que ambos queden mejor que en el caso descentralizado:
\[
\alpha(2450.25)>612.56
\quad\Rightarrow\quad
\alpha>0.25,
\]
\[
(1-\alpha)(2450.25)>1225.13
\quad\Rightarrow\quad
\alpha<0.50.
\]
Por lo tanto, cualquier:
\[
\boxed{0.25<\alpha<0.50}
\]
mejora a ambos agentes respecto de la cadena descentralizada.
\end{formulas}

\begin{ejercicio}{Interpretación de clase: Blockbuster, Netflix y reparto de ingresos}
El caso Blockbuster/Netflix es la aplicación práctica del contrato anterior. Si Blockbuster debía comprar películas a alto costo fijo por unidad, pedía pocas copias. Eso generaba quiebres de stock: clientes querían arrendar un estreno, no lo encontraban y se perdía venta para toda la cadena.

La coordinación consistió en bajar el costo unitario de las copias hacia el costo marginal y compensar a los estudios con una fracción de los ingresos por arriendo. Así el minorista tenía incentivos a pedir más unidades, acercando la cantidad al óptimo centralizado.

En streaming ocurre algo parecido: cuando Netflix paga por visualizaciones o reparte ingresos con estudios, coordina parcialmente la cadena. Cuando produce contenido propio, integra verticalmente y reduce la doble marginalización.
\end{ejercicio}

\begin{tipprueba}{Frase para caso}
``El contrato de reparto de ingresos reduce la doble marginalización porque baja el costo marginal percibido por el retailer y le permite elegir una cantidad más cercana al óptimo integrado. La ganancia total aumenta y luego se reparte con \(\alpha\)''.
\end{tipprueba}

\subsection{Heurísticas V/F y Trampas Clásicas}

\begin{itemize}[itemsep=1ex]
    \item \textbf{Causas Fundamentales (Memorízalas):}
    \begin{enumerate}
        \item \textit{Pronóstico secuencial:} Basar tu producción en lo que te pide tu cliente intermediario en vez del usuario final.
        \item \textit{Order Batching (Loteo):} Acumular pedidos para llenar un camión (distorsiona el ritmo natural).
        \item \textit{Fluctuación de precios (Promociones):} Induce el \textit{Forward Buying} (el cliente compra mucho hoy porque está barato, no porque lo necesite).
        \item \textit{Juegos de escasez:} El cliente infla sus pedidos porque sabe que el proveedor va a racionar el producto.
    \end{enumerate}
    
    \item \textbf{El Juego de la Cerveza (Beer Game):}
    \begin{itemize}
        \item \textit{``El caos del juego se produce porque la demanda final es muy volátil.''} $\rightarrow$ \textbf{FALSO.} En el juego, la demanda es estable. El caos se produce por la \textbf{estructura del sistema} (falta de comunicación y los retrasos físicos/lead times).
    \end{itemize}
    
    \item \textbf{Industria 4.0 vs 5.0:}
    \begin{itemize}
        \item 4.0 = Digitalización pura, IoT, Big Data, hiper-optimización.
        \item 5.0 = Centrada en lo humano, sostenibilidad, \textbf{resiliencia} y \textit{cobots} (humanos + robots).
    \end{itemize}
\end{itemize}
